\section{Main theorem}

\begin{theorem}[Bounded kernel-admission theorem]
For the explicit finite kernel \(K_{900}\) recorded in the repository
\[
\texttt{rookinc/g900-kernel-admission}
\]
at tag
\[
\texttt{g900-kernel-admission-bounded-qed-v1.0.0},
\]
the generated graph \(X_\sigma\) satisfies the following bounded claim set.

\begin{enumerate}[leftmargin=2em]
\item \(X_\sigma\) has 900 vertices and 3600 edges.
\item \(X_\sigma\) is 8-regular, with degree split \(4+4\) into internal and carrier neighbors.
\item The generated edge set equals the recorded canonical edge set.
\item \(X_\sigma\) is connected.
\item \(X_\sigma\) has diameter 8 and radius 6.
\item \(X_\sigma\) separates from the untwisted product baseline by exact metric invariants.
\item The recorded sibling signing, defined in Section 6.1, is not switching-equivalent to the canonical signing.
\item The recorded sibling graph is separated from \(X_\sigma\) by exact graph invariants.
\end{enumerate}

No claim is made here of uniqueness among all signed carriers, external census identity, physical interpretation, or invalidity of the sibling graph.
\end{theorem}

\begin{proof}
The first two items are direct finite consequences of the construction kernel. The remaining items are certified by explicit certificate files generated from the finite payload. These certificates are finite computations, including edge-set comparison, BFS connectedness, all-source BFS metric checks, and an \(\mathbb F_2\) switching-equation rank check; their outputs are included in the repository and summarized in Sections 4--6. The QED status ledger records six direct kernel claims and seven certified claims, with no planned or uncertified claims remaining in the bounded claim set.
\end{proof}

